% Part: many-valued-logic
% Chapter: syntax-and-semantics
% Section: matrices

\documentclass[../../../include/open-logic-section]{subfiles}

\begin{document}

\olfileid{mvl}{syn}{mat}

\olsection{मॅट्रिक्स}

बहुमूल्य तर्कशास्त्र त्याची भाषा, सत्यतामूल्यांचा~$V$ संच,
निर्दिष्ट सत्यतामूल्यांचा उपसंच आणि त्याच्या संयोजकांसाठीची
सत्यता-फलने यांवरून परिभाषित होते. या घटकांना एकत्रितपणे
\emph{मॅट्रिक्स} म्हणतात.

\begin{defn}[मॅट्रिक्स]
\ollabel{defn:matrix}
तर्कशास्त्र~$\Log L$ साठीच्या \emph{मॅट्रिक्स} मध्ये पुढील घटक असतात:
\begin{enumerate}
\item भाषा~$\Lang L$ बनवणारा संयोजकांचा संच;
\item सत्यतामूल्यांचा $V \neq \emptyset$ संच;
\item निर्दिष्ट सत्यतामूल्यांचा $V^+ \subseteq V$ संच;
\item $\Lang L$ मधील प्रत्येक $n$-स्थानी संयोजक~$\star$ साठी
  सत्यता-फलन~$\tf{\star} : V^n \to V$. $n = 0$ असल्यास
  $\tf{\star}$ हे फक्त~$V$ मधील एक मूल्य असते.
\end{enumerate}
\end{defn}

\begin{ex}
अभिजात तर्कशास्त्र~\LogCL{} साठीच्या मॅट्रिक्समध्ये पुढील घटक असतात:
\begin{enumerate}
  \item $\lfalse$, $\lnot$, $\land$, $\lor$, $\lif$ असलेली
  मानक विधानीय भाषा~$\Lang L_0$.
  \item सत्यतामूल्यांचा संच $V = \{\True, \False\}$.
  \item $\True$ हेच एकमेव निर्दिष्ट मूल्य, म्हणजे $V^+ = \{\True\}$.
  \item $\lfalse$ साठी $\tf{\lfalse} = \False$. इतर सत्यता-फलने
  नेहमीच्या सत्यताकोष्टकांनी दिली आहेत (पहा
  \olref{fig:tf-CL}).
\end{enumerate}
\begin{figure}
  \begin{center}
      \begin{tabular}{c|c}
        $\tf{\lnot}$ & \\
        \hline
        $\True$ & $\False$ \\
        $\False$ & $\True$
      \end{tabular}
      \quad
      \begin{tabular}{c|cc}
        $\tf{\land}$ & $\True$ & $\False$ \\
        \hline
        $\True$ & $\True$ & $\False$ \\
        $\False$ & $\False$ & $\False$
      \end{tabular}
      \quad
      \begin{tabular}{c|cc}
        $\tf{\lor}$ & $\True$ & $\False$ \\
        \hline
        $\True$ & $\True$ & $\True$ \\
        $\False$ & $\True$ & $\False$
      \end{tabular}
      \quad
      \begin{tabular}{c|cc}
        $\tf{\lif}$ & $\True$ & $\False$ \\
        \hline
        $\True$ & $\True$ & $\False$ \\
        $\False$ & $\True$ & $\True$
      \end{tabular}
    \end{center}
    \caption{अभिजात तर्कशास्त्र~$\LogCL$ साठीची सत्यता-फलने.}
    \ollabel{fig:tf-CL}
  \end{figure}
\end{ex}

\end{document}
